\documentclass[../../../include/open-logic-section]{subfiles}

\begin{document}

\olfileid[hi]{sth}{replacement}{intro}
\olsection{परिचय}

प्रतिस्थापन वह अभिगृहीत योजना है जो $\ZF$ और~$\Z$ में अंतर करती है।
\crefrange{sth:ordinals::chap}{sth:spine::chap} में हमने पूरे समय उसका
उपयोग किया। इस अध्याय में हम अंततः इस प्रश्न पर विचार करेंगे: क्या
प्रतिस्थापन उचित है?

प्रश्न को तीक्ष्ण बनाने के लिए यह देखना उपयोगी है कि प्रतिस्थापन वास्तव में
बहुत \emph{प्रबल} है। इस अध्याय में (और फिर
\olref[sth][card-arithmetic][fix]{sec} में) हमें आभास होगा कि वह कितना प्रबल
है। इससे संकेत मिलेगा कि उसका औचित्य सचमुच आवश्यक है।

हम दो प्रकार के औचित्य पर चर्चा करेंगे। मोटे तौर पर: \emph{बाह्य} औचित्य
किसी अभिगृहीत को उसके फलों द्वारा उचित ठहराने का प्रयास है; \emph{आंतरिक}
औचित्य यह सुझाकर उसे उचित ठहराने का प्रयास है कि संबंधित गणितीय संकल्पनाएँ
उसका समर्थन करती हैं। इस अध्याय में इसका अर्थ और स्पष्ट होगा, किंतु यह केवल
हिमशैल का सिरा है। अधिक के लिए विशेषतः \citeauthor{Maddy1988a}
(\citeyear{Maddy1988a} और \citeyear{Maddy1988b}) देखें।

\end{document}
